\chapter{Literature Review}
\label{chap:literature-review}

\section{Themes}
\label{sec:literature-themes}

The literature relevant to this thesis is organized around a set of closely related themes concerning the use of operational data to understand, anticipate, and manage airport arrival pressure. These themes include predictive analytics in airport operations, passenger-flow forecasting, baggage-pressure forecasting, congestion and queue management, simulation-based operational evaluation, and decision-support systems. Although these areas are analytically connected, they are often treated as separate research problems. The present review, therefore, examines both the contribution of each theme and the extent to which the literature connects analytical outputs to practical airport management.

\subsection{Predictive analytics and passenger-flow forecasting}
\label{subsec:predictive-analytics-passenger-flow}

Predictive analytics represents a progression from retrospective operational reporting towards anticipatory management. Traditional monitoring systems describe current or historical operating conditions through measures such as flight movements, passenger volumes, processing times, queue lengths, and service performance. These measures are necessary for performance assessment, but by themselves, they do not indicate how operational pressure is likely to develop. Predictive analytics extends monitoring by using temporal and operational information to estimate future demand and identify periods when available capacity may be under increased pressure.

Within airport operations, predictive methods have been applied to flight delays, passenger movements, queue formation, and capacity-related problems. The reviewed material includes time-series models, machine-learning approaches, queueing methods, simulation, optimization, and passenger-tracking techniques \parencite{orsini2019neural,shao2019flight,zhang2021spatiotemporal,shao2022predicting,zhu2024spatiotemporal,pang2024machine}. These approaches differ in their analytical purpose. Forecasting models estimate future demand or operational states; queueing methods examine the relationship between arrivals, service capacity, and waiting; simulation evaluates how processes interact under specified conditions; and optimization supports the selection of resource configurations or operational interventions. Together, these methods provide a basis for moving from descriptive analysis to proactive operational planning.

Passenger-flow forecasting is especially relevant to arrival management because the operational consequences of a flight depend not only on its scheduled or actual arrival time but also on the number of passengers entering the terminal system. When several flights arrive within a short interval, their passenger volumes are combined within the same operational processes. Immigration, baggage reclaim, customs, and terminal exit must therefore manage aggregated demand rather than isolated flights.

\textcite{orsini2019neural} illustrate this problem directly by predicting terminal passenger behavior from anonymized WiFi traces collected at Bologna Airport and by comparing feedforward and recurrent neural network architectures under multi-step forecasting strategies. Their finding that a direct LSTM strategy performed best at short horizons of approximately twenty minutes or less is particularly relevant to arrival-side operations, where staffing and queue-management decisions are only actionable within a similarly narrow window. Temporal characteristics are consequently central to passenger-flow analysis, because pressure varies with the timing and concentration of arrivals rather than with the volume of a single flight considered in isolation. A period with a high number of flights does not necessarily create the same operational conditions as another period unless the passenger loads, timing, and available resources are also considered. Passenger pressure must therefore be interpreted as an operational condition produced by the interaction between arrival demand and processing capacity.

Complementary evidence is provided by studies that model flight-level delay rather than terminal passenger behavior. \textcite{shao2019flight} and the related later study by \textcite{shao2022predicting} represent airport-surface conditions as a situational-awareness map and use engineered traffic-complexity features and a trajectory-based convolutional architecture (TrajCNN) to forecast departure and arrival delay at Los Angeles International Airport, reporting that airport-complexity features are among the principal drivers of predictive performance. \textcite{zhang2021spatiotemporal} extend this approach by combining route-linked weather, airspace, and air-traffic variables within a stacked LSTM architecture, showing that spatio-temporal flight-tracking data of the kind resembled by the FR24-like records used in this thesis can materially improve delay-prediction robustness. \textcite{zhu2024spatiotemporal} further demonstrate, using a network of the busiest airports in China, that delay is not purely a local phenomenon: a self-corrective, causality-informed graph neural network captured inter-airport propagation effects more effectively than simpler distance-based representations. These studies collectively indicate that operational context, temporal persistence, and network effects outperform schedule-based forecasting alone, a conclusion consistent with the temporal and operational feature engineering adopted later in this thesis.

The literature establishes the analytical relevance of passenger and delay forecasting, but it also reveals a distinction between prediction as a technical output and prediction as operational intelligence. A numerical passenger or delay estimate does not automatically indicate the severity of an expected pressure period or the action required. For a forecast to support operations, it must be translated into indicators, pressure levels, congestion classifications, or decision thresholds that correspond to the airport's resource-planning requirements. This distinction is central to the present thesis because the objective is not only to predict arrival pressure but also to convert predictions into manager-facing information.

\subsection{Baggage-pressure forecasting}
\label{subsec:baggage-pressure-forecasting}

Baggage handling forms a distinct but closely related component of the arrival process. Passenger volumes influence baggage demand, but the two pressures are not operationally identical. Passenger pressure refers to the number and concentration of passengers moving through immigration, baggage reclaim, customs, and the terminal exit. Baggage pressure concerns the volume of baggage that must be unloaded, transferred to the reclaim system, delivered to passengers, and processed within the available baggage-handling capacity.

The reviewed literature, spanning passenger-behavior prediction \parencite{orsini2019neural}, flight-delay forecasting \parencite{shao2019flight,shao2022predicting,zhang2021spatiotemporal,zhu2024spatiotemporal}, and arrival-sequencing optimization \parencite{pang2024machine}, places greater emphasis on flight delays and passenger queues than on explicitly forecasting arrival-side baggage demand and baggage-reclaim pressure. Baggage processes are frequently included in broader airport operations, but they are less often treated as an independent predictive problem. This limits airport managers' ability to anticipate periods when baggage-handling resources or reclaim capacity may become constrained, even when passenger-processing resources remain sufficient.

The distinction is operationally important because baggage demand can affect the passenger journey after immigration processing. Passengers may pass through immigration efficiently, but still experience a prolonged arrival process if baggage delivery is delayed. Baggage-reclaim congestion can also affect downstream customs processes and passenger circulation within the reclaim area. Arrival performance must therefore be assessed through both passenger and baggage indicators.

The material underlying this thesis treats baggage pressure as a forecastable operational variable rather than as a secondary consequence of passenger demand. Flight-level arrival information is transformed into baggage-demand indicators that can be aggregated over time and incorporated into forecasting, simulation, and dashboard reporting. This treatment allows distinguishing between periods of passenger pressure, periods of baggage pressure, and periods in which both pressures occur simultaneously.

The literature also indicates that baggage forecasting must be operationally interpretable. An estimate of baggage volume has limited managerial value unless it can be related to baggage-reclaim capacity, expected processing conditions, pressure periods, and possible interventions. The required analytical progression is therefore similar to that of passenger forecasting: operational data are converted into a demand estimate; the estimate is interpreted in relation to capacity; and the resulting pressure information is presented in a form that supports planning.

\subsection{Congestion, queue management, and simulation-based evaluation}
\label{subsec:congestion-queue-simulation}

Congestion occurs when the demand placed on an airport's processes approaches or exceeds the capacity to serve it. In arrival operations, this condition may arise at immigration counters, baggage-reclaim facilities, customs checkpoints, or other connected stages of the passenger journey. The operational effect is not restricted to a single process. Queues formed at one stage may increase passenger system time, alter the rate at which passengers reach subsequent stages, and contribute to wider congestion across the arrival system.

Queue-management research provides a conceptual framework for examining the relationships among passenger arrivals, service resources, waiting, and throughput. \textcite{cao2025robust} demonstrate this relationship directly, applying chance-constrained optimization to airport security-queue management at Barcelona Airport under uncertain passenger compliance and reporting an approximately 85 percent reduction in total waiting time relative to a deterministic baseline. Their finding that explicitly modeling passenger-behavior uncertainty materially improved solution quality reinforces the view, adopted in this thesis, that forecasting and operational response must be linked rather than treated as separate exercises. \textcite{pang2024machine} provide a related demonstration in an arrival-scheduling context: using historical recordings from the Atlanta Air Route Traffic Control Center, they embed a multi-stage machine-learning predictor of separation-time distributions within a mixed-integer linear programming landing-scheduling model, achieving an average 17.2 percent reduction in total landing time compared with first-come-first-served sequencing. These studies show that queueing and optimization methods can support the analysis of individual processes and resource requirements, and that combining prediction with optimization produces larger operational gains than prediction alone.

Airport arrivals, however, involve sequential, interdependent stages. A change to immigration capacity may reduce the immigration queue while increasing the rate at which passengers enter baggage reclaim. Similarly, an improvement in baggage processing may alter the flow towards customs. Evaluation at only one stage may therefore overlook consequences elsewhere in the system. \textcite{sang2026integrating} address this interdependence through an integrated simulation-optimization framework for airport security checkpoints, in which a stochastic optimization model generates delay-aware lane-opening schedules, a mixed-integer linear programming projection ensures workforce-feasible implementation, and an agent-based simulation evaluates operational performance using passenger-arrival data derived from LiDAR sensing rather than static or exogenous arrival profiles. Their approach illustrates the value of representing passenger arrivals as a dynamic, data-driven input to simulation rather than as a fixed assumption, a principle directly relevant to the AnyLogic simulation developed for this thesis.

Simulation provides a means of representing these interactions more broadly. Discrete-event simulation models passengers as they move through a sequence of processes and allows queues, resource use, delays, throughput, and congestion to emerge from the interaction between demand and capacity. In the framework developed for this thesis, AnyLogic simulation connects passenger generation with immigration processing, baggage reclaim, customs inspection, and passenger exit. This structure enables evaluation of the arrival system as an integrated operational process.

Simulation-based evaluation also enables comparison between alternative operating conditions. The simulation scenarios used in the thesis include baseline operations, capacity optimization, peak-stress conditions, dynamic intervention, and integrated optimization. These scenarios allow the consequences of predicted demand and alternative resource decisions to be assessed through common operational indicators. Relevant measures include passenger system time, immigration queue size, baggage queue size, customs queue size, passenger throughput, congestion severity, and operational resilience.

The reviewed literature supports the use of simulation and optimization for scenario testing, but simulation outputs do not automatically constitute operational decisions. A scenario may demonstrate that a particular configuration reduces queues or improves throughput, yet managers still require a clear explanation of the conditions under which the configuration should be applied. The analytical value of simulation, therefore, increases when it is connected to forecasts, congestion thresholds, operational KPIs, and manager-facing recommendations.

The concept of operational resilience extends this evaluation beyond performance under normal conditions. In the context of the present thesis, resilience concerns the ability of the arrival system to maintain or recover acceptable operational performance when demand increases or when pressure is concentrated within peak periods. Simulation enables examination of whether an intervention improves only average performance or also strengthens the system's response under more demanding conditions.

\subsection{Decision-support systems and manager-facing intelligence}
\label{subsec:decision-support-manager-facing}

Decision-support systems provide the interface between analytical outputs and operational decision-making. Forecasting, simulation, and KPI calculation generate information, but airport managers need it organized by operational priority. A decision-support system must therefore make predicted pressure, current performance, scenario outcomes, and recommended actions accessible within a coherent management environment.

The literature reviewed for the thesis indicates that many analytical studies remain focused on model construction and predictive or optimization performance. Even the studies that most closely combine prediction with operational response, the chance-constrained queue-management model of \textcite{cao2025robust}, the ML-enhanced landing-scheduling framework of \textcite{pang2024machine}, and the simulation-optimization checkpoint-scheduling approach of \textcite{sang2026integrating}, report their contributions primarily in terms of waiting-time or landing-time reduction rather than through a manager-facing interface that airport staff across different roles could use directly. Forecast and optimization accuracy are essential, but an accurate model does not by itself ensure that predictions can be used operationally. Managers must be able to identify when pressure is expected, determine which process is affected, assess the likely severity, compare possible responses, and understand the operational implications of each response.

Manager-facing predictive intelligence addresses this requirement by translating technical outputs into interpretable measures. These may include passenger-pressure indicators, baggage-pressure indicators, congestion classifications, KPI summaries, threshold alerts, pressure windows, scenario rankings, and operational recommendations. The purpose is not to remove managerial judgment but to provide structured evidence that supports timely and consistent decisions.

Power BI is used in this thesis as the decision-support environment for presenting forecasting and simulation outputs. The dashboard system combines operational monitoring, predicted pressure, scenario comparison, congestion-risk assessment, dynamic-intervention information, executive recommendations, and explanations of the factors associated with pressure. This integration allows different analytical outputs to be viewed as components of one operational decision process rather than as separate technical results.

The decision-support principle adopted in the thesis follows a progression from data to operational action. Flight-level data are first transformed into passenger and baggage indicators. Forecasting estimates the future development of those indicators. Simulation evaluates the operational consequences of demand and resource configurations. KPIs summarise performance, while the dashboard communicates the results in relation to resource-planning decisions. The DSS therefore forms the final stage of the explain-predict-support logic established in \Cref{chap:introduction}.

This approach is also consistent with the multi-stakeholder character of airport arrival operations. Immigration officers require advance information concerning passenger peaks and staffing pressure. Airline and ground-handling teams require information related to baggage demand and service performance. Airport managers require a consolidated view of the arrival system and its principal bottlenecks. The value of a DSS lies in its ability to organize these different operational concerns without reducing arrival pressure to a single measure.

\section{Gaps}
\label{sec:literature-gaps}

The material reviewed for this thesis identifies five specific and interrelated gaps in the existing literature. These concern the operational focus of previous research, the treatment of baggage demand, the integration of analytical methods, the translation of results into managerial intelligence, and the geographical distribution of airport research.

The first gap is the lack of arrival-centric studies. Airport analytics research often examines flight performance, delays, terminal capacity, security processes, and overall passenger flows. Although these areas are relevant, they do not always examine the arrival system as a distinct sequence of interconnected processes. The delay-prediction studies reviewed in this chapter \parencite{shao2019flight,shao2022predicting,zhang2021spatiotemporal,zhu2024spatiotemporal} and the arrival-scheduling and security-queue studies \parencite{pang2024machine,cao2025robust,sang2026integrating} each concentrate on a single operational target (departure or arrival delay, landing sequencing, or a security checkpoint) rather than on the connected immigration-baggage-customs-exit sequence that defines arrival-side operations at AIBD. Immigration, baggage reclaim, customs, and terminal exit create operational dependencies that differ from those found in departure operations or in a single checkpoint. Research that focuses on only one stage may not capture how pressure propagates throughout the entire arrival journey.

This gap is particularly significant during arrival banks. When several flights converge within short time windows, passenger and baggage demand accumulate across the same resources. The resulting pressure is not adequately represented by flight-level analysis alone. An arrival-centric perspective must account for aggregated demand, sequential processing, queue interactions, baggage delivery, customs flow, and their combined effects on passenger system time.

The second gap is the lack of baggage-pressure forecasting. The reviewed literature provides substantial attention to passenger flows, queues, flight delays, and broader capacity problems. Still, none of the eight studies examined in this chapter explicitly forecasts baggage demand or baggage-reclaim pressure. Baggage is commonly included within general operational descriptions without being modeled as a separate source of arrival pressure.

This omission restricts the operational completeness of passenger-arrival analysis. Passenger and baggage demand are related, but they can create different capacity requirements and bottlenecks. A passenger-flow forecast alone may not identify pressure within baggage handling or reclaim. An integrated arrival analysis, therefore, requires distinct passenger- and baggage-pressure indicators that can be forecast, simulated, and interpreted together.

The third gap is the lack of integrated frameworks combining forecasting, simulation, and decision support. Existing studies commonly focus on a single analytical stage. Forecasting research evaluates predictive performance \parencite{orsini2019neural,shao2019flight,shao2022predicting,zhang2021spatiotemporal,zhu2024spatiotemporal}; queueing and simulation-optimization research evaluates process behavior \parencite{cao2025robust,pang2024machine,sang2026integrating}; but none of the reviewed studies present a dashboard or comparable management interface that communicates its outputs to operational decision-makers. Fewer studies connect prediction, simulation, and dashboard reporting within a continuous operational framework.

The absence of integration creates a separation between anticipated demand and evaluated operational response. Forecasting may indicate that a peak is approaching, but it may not show how that peak will affect queues or system time. Simulation may compare capacity configurations without being connected to predicted short-horizon demand. The literature, therefore, provides limited evidence of frameworks that transform operational data into predictions, simulated consequences, KPI evaluations, and management actions.

The fourth gap is the lack of manager-facing predictive intelligence. Many of the studies reviewed produce technically valid models. Still, their outputs are reported as statistical accuracy, waiting-time reduction, or landing-time reduction rather than as information directly translated for airport managers. Measures such as model errors, predicted values, or simulated queue distributions require operational interpretation before they can support staffing, baggage capacity, or queue management decisions.

Manager-facing intelligence requires more than presenting analytical results. It requires organizing those results around operational questions: when pressure will occur, which process will be affected, how severe the expected condition is, which intervention is appropriate, and what operational effect the intervention is expected to produce. The limited connection between predictive modeling and practical decision-support mechanisms remains a central gap.

The fifth gap is the lack of research on African airports. The eight studies examined in this chapter are situated at Bologna Airport \parencite{orsini2019neural}, Los Angeles International Airport \parencite{shao2019flight,shao2022predicting}, unspecified networked airports studied through ADS-B-style data \parencite{zhang2021spatiotemporal}, the seventy-four busiest airports in China \parencite{zhu2024spatiotemporal}, Atlanta \parencite{pang2024machine}, Barcelona \parencite{cao2025robust}, and airports in the Republic of Korea \parencite{sang2026integrating}. Sub-Saharan African and West African operational contexts receive no attention within this body of work. As a result, the representation of regional airport conditions, institutional coordination requirements, growth patterns, and locally relevant service priorities remains limited.

This geographical gap does not imply that analytical methods developed elsewhere are irrelevant to African airports. It indicates that their application and integration require attention to the specific operational context in which decisions are made. In the case of AIBD, continued traffic growth, concentrated arrival demand, coordination among multiple agencies, the airport's distance from Dakar, and the commitment to a customer-focused form of Senegalese hospitality shape the operational problem addressed by the research.

The five gaps reinforce one another. The limited attention to arrival-side operations contributes to the weak treatment of baggage pressure. The separation of forecasting, simulation, and decision support restricts the operational use of analytical models. The absence of manager-facing intelligence limits the translation of research into practice. The geographical concentration of existing work leaves these issues insufficiently examined in a West African airport environment. The research problem addressed in this thesis is located at the intersection of these gaps.

\section{Positioning}
\label{sec:literature-positioning}

This thesis is positioned within the field of airport operational analytics, with a specific focus on predictive arrival management. It draws on the established analytical foundations of passenger-flow forecasting \parencite{orsini2019neural}, delay and network-effect prediction \parencite{shao2019flight,shao2022predicting,zhang2021spatiotemporal,zhu2024spatiotemporal}, queueing and scheduling optimization \parencite{cao2025robust,pang2024machine}, and simulation-based evaluation \parencite{sang2026integrating}. Its contribution lies in organizing these elements within an arrival-centric predictive operational intelligence framework developed for AIBD.

The study responds to the lack of arrival-centric research by treating the arrival process as an integrated operational system. The unit of interest is not limited to the landing of an aircraft or the performance of an individual checkpoint. Instead, the framework considers the combined movement of passengers and baggage through immigration, baggage reclaim, customs, and terminal exit. This perspective reflects how arrival pressure develops operationally: demand is generated by flights, concentrated within arrival banks, and transmitted through connected processing stages.

The framework explicitly addresses the baggage side of the arrival chain, which the reviewed studies largely omit, by modeling baggage-reclaim capacity as a constraint rather than as an implicit consequence of passenger processing. Passenger and baggage pressures are derived from the validated master dataset of more than 22,000 commercial arrival records and are represented through separate indicators with separate classification rules. The contribution is therefore an operational one rather than a forecasting one: as Section 5.1 establishes, the baggage series used here is derived from the passenger estimate by a fixed multiplier and is not an independent demand signal, so the gap in baggage-demand forecasting identified above is narrowed rather than closed, and closing it would require observed baggage counts of the kind identified in Section 6.4.

The thesis responds to the lack of integrated analytical frameworks by connecting forecasting, simulation, KPI evaluation, and decision support. The framework begins with FR24-like commercial arrival data and the temporal and operational features derived from those data. Machine-learning models are used to estimate short-horizon passenger pressure, baggage demand, and congestion periods, following the short-horizon orientation shown to be effective by \textcite{orsini2019neural}. The resulting predictions are connected to AnyLogic simulation metrics, allowing their operational implications to be evaluated through queue lengths, system time, throughput, congestion severity, and resilience, in a manner consistent with the demand-responsive simulation approach of \textcite{sang2026integrating} and the prediction-linked optimization approach of \textcite{cao2025robust} and \textcite{pang2024machine}.

Simulation scenarios connect predicted pressure to possible operational responses. Baseline conditions establish a reference point, while capacity-optimization, peak-stress, dynamic-intervention, and integrated-optimization scenarios examine how alternative operational configurations perform under varying demand levels. The use of common KPIs enables systematic comparison across these scenarios and supports the interpretation of model outputs in operational terms.

The Power BI Decision Support System completes the integration by presenting forecasts, pressure indicators, congestion thresholds, scenario outcomes, KPIs, and operational recommendations within a manager-facing environment. In this structure, the dashboard is not treated as a separate visualization exercise. It is the operational interface through which analytical and simulation results are translated into information for resource planning, the element identified in \Cref{sec:literature-gaps} as absent from the reviewed literature.

The framework consequently addresses the gap in manager-facing predictive intelligence. It does not stop at generating passenger or baggage forecasts. Predictions are classified and interpreted in relation to operational pressure, incorporated into simulation-based evaluation, and communicated through decision-support outputs. This progression supports decisions concerning staffing, baggage capacity, queue management, and operational intervention.

The study is also positioned in response to the limited representation of African airports in the existing literature. Its empirical and operational setting is Blaise Diagne International Airport in Senegal, a rapidly developing West African airport with ambitions to strengthen its regional hub role. The framework is tailored to AIBD's arrival environment, including the effects of arrival banks, the sequential immigration-baggage-customs flow, coordination across multiple organizations, and the airport's emphasis on customer-focused modernization.

This contextual positioning is not limited to geography. It incorporates the service expectations and stakeholder requirements that shape arrival management at AIBD. Citizens and tourists expect a smooth and welcoming arrival. Business travelers require predictable processing times. Immigration officers require reliable advance information about passenger peaks. Airline and ground-handling teams depend on baggage-delivery performance, operational coordination, and passenger satisfaction. Airport managers require consolidated, interpretable information to plan resources and respond before congestion becomes severe.

The thesis, therefore, treats predictive operational intelligence as both a technical and an organizational capability. Forecasting provides advanced knowledge of expected demand. Simulation evaluates the likely operational consequences of that demand and the performance of alternative responses. The DSS communicates the resulting evidence to operational stakeholders. These components collectively support a shift from reactive congestion management towards proactive, data-driven decision-making.

Academically, the study is grounded in the analytical themes already established in airport forecasting, queue management, simulation, and decision-support research. Its positioning does not depend on introducing an unrelated modeling concept. Instead, it brings established analytical components together around a clearly defined arrival-management problem. The resulting framework is novel within the scope of the reviewed material because it combines an arrival-centric focus, explicit passenger- and baggage-pressure indicators, short-horizon forecasting, simulation-based operational evaluation, and manager-facing Power BI decision support in a West African airport context.

Operationally, the framework is relevant because it links analytical accuracy to resource-planning requirements. Forecasts are evaluated not only as numerical predictions but also as inputs to decisions. Simulation results are interpreted using KPIs that reflect passenger experience and process performance. Dashboard outputs organize these findings according to expected pressure, operational consequences, and possible interventions. The framework, therefore, provides a structured relationship among explanation, prediction, evaluation, and action.

The thesis can consequently be summarised as an integrated response to the five identified gaps. It addresses the lack of arrival-centric studies by modelling the connected arrival process; the limited attention to baggage by treating baggage-reclaim capacity as a distinct operational constraint, while establishing the limits of the derived baggage series as a demand signal; the lack of integrated frameworks by linking forecasting, AnyLogic simulation, KPIs, and Power BI; the lack of manager-facing predictive intelligence by translating outputs into operational indicators and recommendations; and the lack of African airport research by applying the framework to AIBD.

This positioning establishes the basis for the methodology presented in the following chapter. The methodological design operationalizes the literature review by preparing a validated arrival dataset, constructing passenger- and baggage-pressure indicators, engineering temporal and operational features, developing forecasting models, implementing the AnyLogic simulation framework, defining performance indicators, and constructing the Power BI Decision Support System.